% ============================================================
\chapter{SQL --- Requ\^etes de Base}
\label{ch:sql_base}
% ============================================================

En 1974, Donald Chamberlin et Raymond Boyce, chercheurs chez IBM, cr\'eent SEQUEL (\emph{Structured English Query Language}), rebaptis\'e SQL pour des raisons juridiques. Leur ambition~: permettre \`a des non-programmeurs d'interroger des bases de donn\'ees en \'ecrivant des phrases proches de l'anglais. L'id\'ee \'etait r\'evolutionnaire~: au lieu de dire \emph{comment} parcourir les donn\'ees, on dit simplement \emph{ce que} l'on veut. Cinquante ans plus tard, SQL reste le langage universel des bases de donn\'ees relationnelles.

\section{Introduction \`a SQL}

\begin{definition}[SQL]
\emph{SQL} (\emph{Structured Query Language}) est le langage standard
pour interroger et manipuler les bases de donn\'ees relationnelles.
Contrairement \`a l'alg\`ebre relationnelle, SQL est un langage
\emph{d\'eclaratif}~: on sp\'ecifie \emph{ce que} l'on veut, pas \emph{comment}
l'obtenir.
\end{definition}

SQL a \'et\'e standardis\'e par l'ISO/ANSI. Les versions majeures sont~:
SQL-86, SQL-92, SQL:1999, SQL:2003, SQL:2011, SQL:2016, SQL:2023.

\section{Structure d'une requ\^ete SELECT}

\begin{formules}[title=\textbf{Syntaxe g\'en\'erale du SELECT}]
\begin{minted}{sql}
SELECT [DISTINCT] colonnes
FROM   tables
[WHERE  condition]
[ORDER BY colonne [ASC|DESC]]
[LIMIT  n [OFFSET m]];
\end{minted}
\end{formules}

L'ordre d'ex\'ecution logique est diff\'erent de l'ordre d'\'ecriture~:
\begin{enumerate}
    \item \texttt{FROM} --- d\'eterminer les tables sources.
    \item \texttt{WHERE} --- filtrer les lignes.
    \item \texttt{SELECT} --- projeter les colonnes.
    \item \texttt{DISTINCT} --- \'eliminer les doublons.
    \item \texttt{ORDER BY} --- trier le r\'esultat.
    \item \texttt{LIMIT / OFFSET} --- limiter le nombre de lignes.
\end{enumerate}

\section{SELECT et FROM}

\begin{codeblock}[title=\textbf{S\'electionner toutes les colonnes}]
\begin{minted}{sql}
SELECT * FROM Etudiants;
\end{minted}
\end{codeblock}

\begin{codeoutput}
\begin{verbatim}
id_etudiant | nom     | prenom | date_naissance | filiere
------------+---------+--------+----------------+--------
1           | Bernard | Alice  | 2004-03-15     | Info
2           | Petit   | Bob    | 2003-11-22     | Info
3           | Moreau  | Claire | 2004-07-08     | Maths
4           | Leroy   | David  | 2003-01-30     | Info
5           | Roux    | Emma   | 2004-09-12     | Maths
\end{verbatim}
\end{codeoutput}

\begin{codeblock}[title=\textbf{S\'electionner des colonnes sp\'ecifiques}]
\begin{minted}{sql}
SELECT nom, prenom, filiere FROM Etudiants;
\end{minted}
\end{codeblock}

\begin{bonnepratique}[title=\textbf{\'Eviter SELECT *}]
En production, \'evitez \texttt{SELECT *}~:
\begin{itemize}
    \item Cela transf\`ere des donn\'ees inutiles.
    \item Le r\'esultat d\'epend du sch\'ema~: l'ajout d'une colonne peut casser le code.
    \item Utilisez \texttt{SELECT *} uniquement pour l'exploration et le d\'ebogage.
\end{itemize}
\end{bonnepratique}

\section{Alias de colonnes et de tables}

\begin{codeblock}[title=\textbf{Utilisation d'alias}]
\begin{minted}{sql}
SELECT e.nom AS nom_etudiant,
       e.prenom AS prenom_etudiant,
       e.filiere AS departement
FROM Etudiants AS e;
\end{minted}
\end{codeblock}

Le mot-cl\'e \texttt{AS} est optionnel dans la plupart des SGBD~:
\texttt{SELECT nom nom\_etudiant FROM Etudiants e}.

\section{Clause WHERE --- Filtrage}

\begin{definition}[Clause WHERE]
La clause \texttt{WHERE} filtre les lignes selon une condition bool\'eenne.
Seules les lignes pour lesquelles la condition est \texttt{TRUE} sont retenues
(\texttt{FALSE} et \texttt{UNKNOWN} sont rejet\'es).
\end{definition}

\subsection{Op\'erateurs de comparaison}

\begin{formules}[title=\textbf{Op\'erateurs de comparaison SQL}]
\begin{center}
\begin{tabular}{ll}
\toprule
\textbf{Op\'erateur} & \textbf{Signification} \\
\midrule
\texttt{=} & \'Egal \\
\texttt{<>} ou \texttt{!=} & Diff\'erent \\
\texttt{<}, \texttt{>} & Inf\'erieur, sup\'erieur \\
\texttt{<=}, \texttt{>=} & Inf\'erieur ou \'egal, sup\'erieur ou \'egal \\
\texttt{BETWEEN a AND b} & Entre $a$ et $b$ (inclus) \\
\texttt{IN (v1, v2, ...)} & Appartient \`a une liste \\
\texttt{LIKE 'motif'} & Correspondance de motif \\
\texttt{IS NULL} / \texttt{IS NOT NULL} & Test de nullit\'e \\
\bottomrule
\end{tabular}
\end{center}
\end{formules}

\begin{codeblock}[title=\textbf{Exemples de filtrage}]
\begin{minted}{sql}
-- Etudiants en Info
SELECT * FROM Etudiants WHERE filiere = 'Info';

-- Notes entre 10 et 15
SELECT * FROM Inscriptions WHERE note BETWEEN 10 AND 15;

-- Etudiants dont le nom commence par 'B'
SELECT * FROM Etudiants WHERE nom LIKE 'B%';

-- Etudiants en Info ou Maths
SELECT * FROM Etudiants WHERE filiere IN ('Info', 'Maths');

-- Inscriptions sans note
SELECT * FROM Inscriptions WHERE note IS NULL;
\end{minted}
\end{codeblock}

\subsection{Op\'erateurs logiques}

\begin{codeblock}[title=\textbf{Combinaison de conditions}]
\begin{minted}{sql}
-- AND : les deux conditions doivent etre vraies
SELECT * FROM Etudiants
WHERE filiere = 'Info' AND date_naissance > '2004-01-01';

-- OR : au moins une condition vraie
SELECT * FROM Inscriptions
WHERE note > 16 OR note < 8;

-- NOT : negation
SELECT * FROM Etudiants
WHERE NOT filiere = 'Maths';
\end{minted}
\end{codeblock}

\begin{attention}[title=\textbf{Logique ternaire et NULL}]
En SQL, la logique est \emph{ternaire}~: \texttt{TRUE}, \texttt{FALSE}, \texttt{UNKNOWN}.
\begin{center}
\begin{tabular}{c|ccc}
\texttt{AND} & T & F & U \\
\hline
T & T & F & U \\
F & F & F & F \\
U & U & F & U \\
\end{tabular}
\qquad
\begin{tabular}{c|ccc}
\texttt{OR} & T & F & U \\
\hline
T & T & T & T \\
F & T & F & U \\
U & T & U & U \\
\end{tabular}
\end{center}
\texttt{NOT UNKNOWN = UNKNOWN}. Cons\'equence~: \texttt{note = NULL}
retourne \texttt{UNKNOWN}, pas \texttt{TRUE}. Utilisez \texttt{IS NULL}.
\end{attention}

\subsection{LIKE et les motifs}

\begin{formules}[title=\textbf{Caract\`eres sp\'eciaux de LIKE}]
\begin{center}
\begin{tabular}{ll}
\toprule
\textbf{Caract\`ere} & \textbf{Signification} \\
\midrule
\texttt{\%} & Z\'ero ou plusieurs caract\`eres quelconques \\
\texttt{\_} & Exactement un caract\`ere quelconque \\
\bottomrule
\end{tabular}
\end{center}
\end{formules}

\begin{codeblock}[title=\textbf{Exemples avec LIKE}]
\begin{minted}{sql}
-- Noms commencant par 'M'
SELECT * FROM Etudiants WHERE nom LIKE 'M%';

-- Prenoms de exactement 5 lettres
SELECT * FROM Etudiants WHERE prenom LIKE '_____';

-- Noms contenant 'or'
SELECT * FROM Etudiants WHERE nom LIKE '%or%';
\end{minted}
\end{codeblock}

\section{ORDER BY --- Tri}

\begin{codeblock}[title=\textbf{Tri des r\'esultats}]
\begin{minted}{sql}
-- Tri ascendant (par defaut)
SELECT nom, prenom FROM Etudiants ORDER BY nom ASC;

-- Tri descendant
SELECT * FROM Inscriptions ORDER BY note DESC;

-- Tri multi-colonnes
SELECT * FROM Etudiants ORDER BY filiere ASC, nom ASC;
\end{minted}
\end{codeblock}

\begin{codeoutput}
\begin{verbatim}
-- Tri par note descendante :
id_etudiant | id_cours | note | annee
-----------+---------+------+------
5           | 104      | 18.0 | 2025
3           | 102      | 17.5 | 2025
2           | 103      | 16.0 | 2025
1           | 101      | 15.5 | 2025
3           | 104      | 14.0 | 2025
1           | 102      | 13.0 | 2025
2           | 101      | 12.0 | 2025
4           | 101      | 9.0  | 2025
\end{verbatim}
\end{codeoutput}

\section{DISTINCT --- \'Elimination des doublons}

\begin{codeblock}[title=\textbf{Utilisation de DISTINCT}]
\begin{minted}{sql}
-- Sans DISTINCT : doublons possibles
SELECT filiere FROM Etudiants;
-- Info, Info, Maths, Info, Maths

-- Avec DISTINCT : valeurs uniques
SELECT DISTINCT filiere FROM Etudiants;
-- Info, Maths
\end{minted}
\end{codeblock}

\section{LIMIT et OFFSET --- Pagination}

\begin{codeblock}[title=\textbf{Limiter les r\'esultats}]
\begin{minted}{sql}
-- Les 3 meilleures notes
SELECT e.nom, i.note
FROM Etudiants e
JOIN Inscriptions i ON e.id_etudiant = i.id_etudiant
ORDER BY i.note DESC
LIMIT 3;

-- Pagination : page 2 (lignes 4 a 6)
SELECT * FROM Etudiants ORDER BY id_etudiant LIMIT 3 OFFSET 3;
\end{minted}
\end{codeblock}

\begin{codeoutput}
\begin{verbatim}
-- Les 3 meilleures notes :
nom   | note
------+-----
Roux  | 18.0
Moreau| 17.5
Petit | 16.0
\end{verbatim}
\end{codeoutput}

\begin{attention}[title=\textbf{LIMIT n'est pas standard SQL}]
\texttt{LIMIT} est support\'e par PostgreSQL, MySQL, SQLite.
Le standard SQL:2008 utilise \texttt{FETCH FIRST n ROWS ONLY}~:
\begin{minted}{sql}
SELECT * FROM Etudiants
ORDER BY nom
FETCH FIRST 3 ROWS ONLY;  -- Standard SQL
\end{minted}
\end{attention}

\section{Expressions et fonctions scalaires}

\begin{codeblock}[title=\textbf{Expressions calcul\'ees}]
\begin{minted}{sql}
-- Expression arithmetique
SELECT nom, note, note / 20.0 * 100 AS pourcentage
FROM Inscriptions i
JOIN Etudiants e ON i.id_etudiant = e.id_etudiant;

-- Concatenation de chaines
SELECT nom || ' ' || prenom AS nom_complet
FROM Etudiants;

-- Fonction conditionnelle CASE
SELECT nom, note,
    CASE
        WHEN note >= 16 THEN 'Tres bien'
        WHEN note >= 14 THEN 'Bien'
        WHEN note >= 12 THEN 'Assez bien'
        WHEN note >= 10 THEN 'Passable'
        ELSE 'Insuffisant'
    END AS mention
FROM Inscriptions i
JOIN Etudiants e ON i.id_etudiant = e.id_etudiant;
\end{minted}
\end{codeblock}

\begin{codeoutput}
\begin{verbatim}
nom     | note | mention
--------+------+-----------
Bernard | 15.5 | Bien
Bernard | 13.0 | Assez bien
Petit   | 12.0 | Assez bien
Petit   | 16.0 | Tres bien
Moreau  | 17.5 | Tres bien
Moreau  | 14.0 | Bien
Leroy   | 9.0  | Insuffisant
Roux    | 18.0 | Tres bien
\end{verbatim}
\end{codeoutput}

\section{Fonctions utiles}

\begin{formules}[title=\textbf{Fonctions scalaires courantes}]
\begin{center}
\begin{tabular}{lll}
\toprule
\textbf{Cat\'egorie} & \textbf{Fonction} & \textbf{Description} \\
\midrule
Texte & \texttt{UPPER(s)}, \texttt{LOWER(s)} & Majuscules, minuscules \\
      & \texttt{LENGTH(s)} & Longueur \\
      & \texttt{SUBSTR(s, d, n)} & Sous-cha\^ine \\
      & \texttt{TRIM(s)} & Suppression espaces \\
      & \texttt{REPLACE(s, a, b)} & Remplacement \\
Nombre & \texttt{ABS(x)}, \texttt{ROUND(x, n)} & Valeur absolue, arrondi \\
       & \texttt{CEIL(x)}, \texttt{FLOOR(x)} & Plafond, plancher \\
Date   & \texttt{CURRENT\_DATE} & Date courante \\
       & \texttt{DATE('now')} & Date courante (SQLite) \\
NULL   & \texttt{COALESCE(a, b, c)} & Premi\`ere valeur non NULL \\
       & \texttt{NULLIF(a, b)} & NULL si $a = b$ \\
\bottomrule
\end{tabular}
\end{center}
\end{formules}

\begin{codeblock}[title=\textbf{Exemples de fonctions}]
\begin{minted}{sql}
SELECT UPPER(nom) AS nom_maj,
       LENGTH(prenom) AS lg_prenom,
       COALESCE(note, 0) AS note_ou_zero
FROM Etudiants e
LEFT JOIN Inscriptions i ON e.id_etudiant = i.id_etudiant;
\end{minted}
\end{codeblock}

\section{INSERT, UPDATE, DELETE}

\begin{codeblock}[title=\textbf{Insertion de donn\'ees}]
\begin{minted}{sql}
-- Insertion d'une ligne
INSERT INTO Etudiants (id_etudiant, nom, prenom, filiere)
VALUES (6, 'Garcia', 'Lucie', 'Info');

-- Insertion multiple
INSERT INTO Etudiants (id_etudiant, nom, prenom, filiere) VALUES
    (7, 'Dubois', 'Marc', 'Maths'),
    (8, 'Thomas', 'Julie', 'Info');

-- Insertion depuis une requete
INSERT INTO Archive_Etudiants
SELECT * FROM Etudiants WHERE filiere = 'Maths';
\end{minted}
\end{codeblock}

\begin{codeblock}[title=\textbf{Mise \`a jour et suppression}]
\begin{minted}{sql}
-- Mise a jour
UPDATE Inscriptions SET note = 11.0
WHERE id_etudiant = 4 AND id_cours = 101;

-- Suppression conditionnelle
DELETE FROM Inscriptions WHERE note < 8;

-- Suppression de toutes les lignes
DELETE FROM Archive_Etudiants;
-- ou TRUNCATE TABLE Archive_Etudiants;  (PostgreSQL)
\end{minted}
\end{codeblock}

\begin{attention}[title=\textbf{UPDATE/DELETE sans WHERE}]
Un \texttt{UPDATE} ou \texttt{DELETE} sans clause \texttt{WHERE} affecte
\emph{toutes} les lignes de la table. Toujours v\'erifier la clause
\texttt{WHERE} avant d'ex\'ecuter~!
\end{attention}

\section{Int\'egration Python compl\`ete}

\begin{codeblock}[title=\textbf{CRUD complet avec Python/sqlite3}]
\begin{minted}{python}
import sqlite3

conn = sqlite3.connect(':memory:')  # base en memoire
cur = conn.cursor()

# Creation et insertion
cur.executescript('''
    CREATE TABLE Etudiants (
        id INTEGER PRIMARY KEY, nom TEXT, prenom TEXT, filiere TEXT
    );
    INSERT INTO Etudiants VALUES (1, 'Bernard', 'Alice', 'Info');
    INSERT INTO Etudiants VALUES (2, 'Petit', 'Bob', 'Info');
    INSERT INTO Etudiants VALUES (3, 'Moreau', 'Claire', 'Maths');
''')

# Lecture avec parametres
filiere = 'Info'
cur.execute("SELECT * FROM Etudiants WHERE filiere = ?", (filiere,))
print("Etudiants Info :", cur.fetchall())

# Mise a jour
cur.execute("UPDATE Etudiants SET filiere = ? WHERE id = ?", ('Maths', 2))
conn.commit()

# Verification
cur.execute("SELECT * FROM Etudiants ORDER BY id")
for row in cur.fetchall():
    print(row)

conn.close()
\end{minted}
\end{codeblock}

\begin{codeoutput}
\begin{verbatim}
Etudiants Info : [(1, 'Bernard', 'Alice', 'Info'), (2, 'Petit', 'Bob', 'Info')]
(1, 'Bernard', 'Alice', 'Info')
(2, 'Petit', 'Bob', 'Maths')
(3, 'Moreau', 'Claire', 'Maths')
\end{verbatim}
\end{codeoutput}

\section*{Exercices}

\begin{exercice}
\'Ecrivez les requ\^etes SQL suivantes~:
\begin{enumerate}
    \item Tous les cours ayant plus de 3 cr\'edits.
    \item Les noms et pr\'enoms des \'etudiants n\'es avant 2004.
    \item Les inscriptions avec une note sup\'erieure \`a 15, tri\'ees par note d\'ecroissante.
    \item Les 3 \'etudiants ayant les meilleures notes.
\end{enumerate}
\end{exercice}

\begin{exercice}
Expliquez la diff\'erence entre les deux requ\^etes~:
\begin{minted}{sql}
SELECT DISTINCT filiere FROM Etudiants;
SELECT filiere FROM Etudiants GROUP BY filiere;
\end{minted}
\end{exercice}

\begin{exercice}
\'Ecrivez une requ\^ete SQL qui affiche pour chaque inscription~: le nom
de l'\'etudiant, l'intitul\'e du cours, la note, et la mention correspondante
(Tr\`es bien $\geq$ 16, Bien $\geq$ 14, Assez bien $\geq$ 12, Passable $\geq$ 10,
Insuffisant sinon).
\end{exercice}

\begin{exercice}
Quels sont les r\'esultats des expressions suivantes ?
\begin{enumerate}
    \item \texttt{NULL = NULL}
    \item \texttt{NULL <> 5}
    \item \texttt{NULL AND TRUE}
    \item \texttt{NULL OR TRUE}
    \item \texttt{NOT NULL}
\end{enumerate}
\end{exercice}

\begin{exercice}
\'Ecrivez un programme Python qui~:
\begin{enumerate}
    \item Cr\'ee une base de donn\'ees SQLite avec les tables du sch\'ema Universit\'e.
    \item Ins\`ere les donn\'ees d'exemple.
    \item Ex\'ecute 3 requ\^etes diff\'erentes et affiche les r\'esultats format\'es.
    \item Utilise des requ\^etes param\'etr\'ees pour toutes les entr\'ees utilisateur.
\end{enumerate}
\end{exercice}
